% Official source: https://www.imo-official.org/assets/documents/problems/2000/2000_eng.pdf
\competitionyear{2000}

\begin{problem}{2000}{1}
Line $AB$ is tangent to the circles through $C, A, M, N$ and $N, M, B, D$. The point $M$ lies between $C$ and $D$, with $\overleftrightarrow{CD}\parallel\overleftrightarrow{AB}$. Chords $NA$ and $CM$ intersect at $P$, and chords $NB$ and $MD$ at $Q$. The rays $\overrightarrow{CA}$ and $\overrightarrow{DB}$ meet at $E$. Prove that $PE=QE$.
\end{problem}

\begin{problem}{2000}{2}
For positive real numbers $A, B, C$ with $ABC=1$, prove that
\[
    \left(A-1+\frac1B\right)\left(B-1+\frac1C\right)\left(C-1+\frac1A\right)\leq1.
\]
\end{problem}

\begin{problem}{2000}{3}
Let $k>0$ be real and $N>1$ an integer. Place $N$ points on a line, not all coincident. In a move, select two points $A, B$ at distinct locations with $A$ to the right of $B$, and replace $B$ by a point $B'$ to the right of $A$ satisfying $AB'=kBA$. Determine the values of $k$ for which repeated moves can move the points arbitrarily far to the right.
\end{problem}

\begin{problem}{2000}{4}
Distribute $100$ cards numbered $1$ through $100$ among three nonempty boxes. Count the distributions for which the sum of the numbers on any two cards chosen from different boxes uniquely determines the third box.
\end{problem}

\begin{problem}{2000}{5}
Does there exist a positive integer $N$ having exactly $2000$ distinct prime divisors and satisfying $N\mid 2^N+1$? Prime factors may occur with multiplicity.
\end{problem}

\begin{problem}{2000}{6}
In acute triangle $A_1A_2A_3$, let $K_i$ be the foot of the altitude from $A_i$, and let $L_i$ be the point where the incircle touches the side opposite $A_i$. Reflect lines $K_1K_2$, $K_2K_3$, and $K_3K_1$ in lines $L_1L_2$, $L_2L_3$, and $L_3L_1$, respectively. Prove that the three reflected lines form a triangle whose vertices lie on the incircle.
\end{problem}
